% Gap Series Part 4; corrective working paper v0.3, 2026-10-02.
% Historical source-bound form audit retained; R19 inventory delta only.
% Original v0.2 files and their claim/source bindings remain protected.
\documentclass[11pt,a4paper]{article}
\usepackage{cmap}
\usepackage[utf8]{inputenc}
\usepackage[T1]{fontenc}
\usepackage{lmodern}
\usepackage{microtype}
\usepackage{amsmath,amssymb,amsthm}
\usepackage{booktabs}
\usepackage{tabularx}
\usepackage{array}
\usepackage[colorlinks=true,linkcolor=blue,citecolor=blue,urlcolor=blue]{hyperref}
\renewcommand*{\HyperDestNameFilter}[1]{en.#1}
\usepackage{xurl}
\input{glyphtounicode}
\pdfgentounicode=1

\hypersetup{
  pdftitle={The Bridge Certificate Applied to the Joshi ATS Series: A Ledger Audit},
  pdfauthor={Lukas Geiger},
  pdfsubject={Pre-specified four-row ledger audit of Kirti Joshi's Arithmetic Teichmüller Space series},
  pdfkeywords={abc conjecture, Inter-universal Teichmüller Theory, Arithmetic Teichmüller Spaces, Corollary 3.12, Bridge Certificate, Ledger Audit},
  pdflang={en},
  bookmarksnumbered=true,
  pdfpagemode=UseOutlines
}

\clubpenalty=10000
\widowpenalty=10000
\displaywidowpenalty=10000
\frenchspacing

\newtheorem{remark}{Remark}[section]

\title{The Bridge Certificate Applied to the Joshi ATS Series:\\A Ledger Audit}
\author{Lukas Geiger\\\small Independent Researcher, Bernau, Germany\\\small ORCID: 0009-0005-7296-1534}
\date{2 October 2026 --- Version 0.3 (working paper)}

\begin{document}
\maketitle

\begin{abstract}
We revisit the pre-specified form audit of Kirti Joshi's Arithmetic
Teichmüller Space (ATS) series, with source versions and the historical
selection of rows R2, R4, R6, R8 preserved. This is a selected-rows form
audit, not a complete certificate or a proof-correctness audit. The source
makes the $j^2$-scaling law, the assigned compensating site, and the
asymmetric Frobenius shift visible. These historical findings have narrower
predicates than the current seven-component contract: locating a load is
not a quantified budget, and seeing a second mode is not a complete E8
mode audit. The revised crosswalk assigns F-R2a to C4 (dependent on C6),
F-R2b and F-R6a to C6, and F-R6b to C7; F-R8 belongs to the triggered
extension E8 outside that crosswalk. A declared normalization convention
is classified as partially delivering (I), rather than as a derivation of
cross-container compatibility. We correct our previous generic assertion
that a set's volume dominates the absolute value of each element: the cited
ATS III argument concerns a tensor-product box under specific integral
and weight hypotheses. A simple translated-ball example refutes only our
generic paraphrase. ATS III box containment and ATS IV's normalization
and measure bridge remain separate obligations. Later normalization and
effective-weight classifications are cited within their stated hypotheses;
they do not close indexed log-volume, hull, measure, or mode compatibility.
No claim about the truth of abc, the correctness of IUT, or the correctness
of the ATS series follows from this audit.
\end{abstract}

\clearpage
\tableofcontents
\clearpage

%=====================================================================
\section{Mandate, Source Versions, and Scope}\label{sec:mandate}

This paper is Part~4 of the abc Gap Series. Part~1~\cite{part1}
analyzed the comparison in the transport step of Corollary~3.12
of IUTchIII~\cite{mochizukiIUT}. Part~2~\cite{part2} defines the
\emph{bridge contract}; Part~3~\cite{part3} supplies its line-level
ledger instrument. Both are published working papers. The present
paper applies selected historical predicates to the exact ATS
versions~\cite{joshiI,joshiII,joshiIIhalf,joshiIII,joshiIV}.
The correction dated 2~October 2026 processes the inventoried
mapping, source-paraphrase, and status-import errors in our v0.2 text;
it performs no new whole-series proof audit.

\subsection{Ground rules}

First, \textbf{no authority argument}: neither ``Joshi asserts it holds'' nor
``Scholze--Stix~\cite{scholzestix} assert it does not'' counts as evidence;
every verdict is anchored in the verbatim wording of the audited sources.
Joshi's own status tables about the Scholze--Stix objections are treated as
self-report, not as evidence. Second, \textbf{no truth verdict}: the audit
judges whether the text serves pre-specified contract clauses, not whether
abc is true, whether Joshi's proof is correct, or whether the critics are
right. Third, \textbf{status discipline}: the audited texts are publicly
posted, versioned arXiv preprints, not peer-reviewed publications; ATS~IV
carries the header ``Preliminary version for comments'' and opens with
``[This is still a work in progress!]''. We therefore write ``publicly posted
text'' throughout, bind exact versions (Table~\ref{tab:manifest}), and treat
the preprint status as a version-stability caveat recorded per finding.

\subsection{Pre-specified scope: a four-row test}\label{sec:scope}

The audit mandate --- fixed in the series' ledger instrument before this test
was run --- restricts the test to four ledger rows. \textbf{This is a
commissioned four-row test, not a complete F-R0--F-R7 core audit
with triggered E8.} Rows R0, R1, R3, R5,
R7 were not part of the mandate (R5 is touched only insofar as R4 forces it).
Table~\ref{tab:audit-scope} fixes the included and excluded rows.

\begin{table}[htbp]
\centering\small
\begin{tabularx}{\textwidth}{@{}l >{\raggedright\arraybackslash}X l@{}}
\toprule
Row & Reason for pre-selection & Not covered \\
\midrule
R2 & numerical pilot transport; comparative discrimination case & no R0/R1 audit \\
R4 & visible treatment of the $j^2$-load & R5 only via R4 \\
R6 & contract components 6--7; the central bridge question & no proof-correctness audit \\
F-R8 & visibility of the second mode & no completeness of all auxiliary modes \\
\bottomrule
\end{tabularx}
\caption{Pre-specified four-row audit scope. The table documents the mandate
boundary and prevents the row findings from being read as a complete
certificate audit.}
\label{tab:audit-scope}
\end{table}

\subsection{Operational verdict rules}\label{sec:rules}

Each atomic claim receives one of three verdicts. \emph{Delivers} $=$ the
pre-stated minimal predicate of the row is served by the quoted text.
\emph{Partially delivers} $=$ at least one explicit sub-predicate remains
unserved. \emph{Does not deliver} $=$ the minimal predicate is unserved.
\textbf{A row-level pass is not a certificate pass}: it certifies only the
row's minimal predicate, never the full seven-component contract. Confidence
is reported in three separate dimensions: \textbf{Q} (source binding: is the
verbatim anchor unambiguous?), \textbf{K} (criterion fit: does the wording
serve the pre-stated predicate?), \textbf{S} (substance status: checked /
not checked / externally marked open).

Operational glossary (minimal pairs used below):
\emph{transported} $=$ a numerical quantity together with an explicit
transformation law; \emph{invariant} $=$ same output value after typed
identification; \emph{addressed} $=$ load and its compensating site named;
\emph{quantitatively booked} $=$ amount/rule balanced in a complete budget;
\emph{visible mode} $=$ mode, role, and affected bound expressly named.

\subsection{Mapping to the current contract}\label{sec:crosswalk}

The historical mandate remains R2/R4/R6 plus R8. In the current
namespace of Part~2~\cite{part2}, the atomic findings below use
F-R2a, F-R2b, F-R4a, F-R4b, F-R6a, F-R6b, and F-R8.
F-R2a belongs to C4 with a dependency on C6; F-R2b and F-R6a to C6;
F-R4a and F-R4b to C5; F-R6b to C7. The typed evaluation identity
$\nu_C\circ i_q=\nu_q$ is a C6/F-R6a obligation, not component~1
and not certified merely by displaying a scaling law.
F-R8 is the triggered E8 extension, outside the seven-component
crosswalk. It is X (not included in a denominator) when no relevant
multi-mode, sign, or ordering trigger occurs; the explicit Frobenius
shift triggers it here.

The historical address/visibility predicates are retained as such.
A historical \emph{delivers} finding for F-R4a or F-R6b does not assert
the current C5 budget or the current C7 hypothesis-by-hypothesis
coverage. Likewise, F-R8 visibility does not establish complete
E8 mode coverage. We therefore report no current full-certificate
score and no ratio of successful rows. Under the current grammar,
a documented convention is at most I (partially delivers), not
E+ (delivers a compatibility equation). Joshi is a comparative
discrimination case; it is not the accepted control of the
separate historical v1-T illustration in Part~2.

\subsection{Source manifest}\label{sec:manifest}

\begin{table}[htbp]
\centering\small
\begin{tabularx}{\textwidth}{@{}l >{\raggedright\arraybackslash}X l@{}}
\toprule
ID & Source (exact version audited) & Role \\
\midrule
ATS~III & arXiv:2401.13508v4 (rev.\ 24~Feb 2025)~\cite{joshiIII} & Cor.~3.12 construction (R2, R4, R6) \\
ATS~IV & arXiv:2403.10430v2 (rev.\ 24~Feb 2025)~\cite{joshiIV} & abc conclusion, bridge site (R2, R6, F-R8) \\
ATS~II($\tfrac12$) & arXiv:2305.10398v12~\cite{joshiIIhalf} & normalization data (status notice A1) \\
\bottomrule
\end{tabularx}
\caption{Audited sources. Line anchors ``extract ll.'' refer to the plain-text
extracts of these exact PDF versions retained with the historical
audit records. They are local finding aids, not independently published
evidence; the exact-version section/theorem citations are the
public primary-source anchors. Audit date: 12~August
2026; status notices dated 13/15~August 2026 (Section~\ref{sec:notices}).}
\label{tab:manifest}
\end{table}

\subsection{\texorpdfstring{Four ``tautology'' passages, disambiguated}{Four "tautology" passages, disambiguated}}\label{sec:tautology}

The audited texts use the word ``tautology'' in four places with distinct
scopes, and conflating them produces the false impression that Joshi calls his
own arithmetic proof a tautology --- which he does not. (T-a) a
measure-theoretic triviality used as a tool (ATS~III, Lemma~9.10.2.2(4));
(T-b) a statement about the \emph{geometric} analog only (ATS~III, \S\,12.18);
(T-c) a statement about Joshi's own variant $\tilde\Theta_{\mathrm{Joshi}}$,
while the theorem in question is proven for
$\tilde\Theta_{\mathrm{Mochizuki}}$ (ATS~IV, \S\,1.6(4e)); and (T-d) the
load-bearing bridge equality in the proof of ATS~IV, Theorem~6.10.1.
\textbf{Only T-d is load-bearing for this audit.}

%=====================================================================
\section{The Ledger Audit (Atomic Findings)}\label{sec:ledger}

\subsection{F-R2a --- intra-link transport: delivers}\label{sec:r2a}

\emph{Minimal predicate: the numerical $q$-readout is transported along the
theta-link by an explicit numerical law, with the $j^2$-factor carried
visibly.} ATS~III, Theorem~4.2.2.1(4) states the valuation scaling relation
$|-|_{K'_{y_{w,j}}} = |-|^{j^2}_{K'_{y_{w,1}}}$ at the odd semistable places
(``Then one has the following relationship between their valuations,''
eq.~(4.2.2.2); extract ll.~1544--1560); \S\,9.9 carries the $q$-value through it
explicitly (eqs.~(9.9.3)--(9.9.4); extract ll.~6458--6484); and \S\,9.10.2
supplies the concrete log-volume scale (normalized by
$\mathrm{LogVol}(O_{L'_w}) = 0$, $\mathrm{Vol}(\lambda O) = |\lambda|$;
extract ll.~6562--6596). This is \emph{transport by an explicit numerical
law}, not invariance of a value: what is certified is a covariant
transformation rule with a concrete numerical target scale, serving the historical F-R2a/C4 transport-law predicate (dependent on C6). The full typed identity $\nu_C \circ i_q = \nu_q$ is not written
out in the source and is not certified here.
\textbf{Verdict: delivers (row predicate). Q high / K high / S not checked.}

\subsection{F-R2b --- cross-container transport: set by convention}\label{sec:r2b}

\emph{Minimal predicate: the equality of readouts across the global Frobenius
shift $y_0 \leftrightarrow \varphi(y_0)$ is derived, not set.} Here the
equality is \emph{set} by the choice of ``the same global normalization''
(ATS~IV, Thm.~6.10.1 preamble; extract ll.~3186--3189), with a one-sentence
proof (``The equality \dots is a consequence of the choice of the same global
normalization \dots''; extract ll.~3262--3266), while the local constituents
change under $\varphi$ and the source itself notes that this ``precludes
direct comparison of various local quantities'' (ATS~IV, Rem.~6.10.2; extract
ll.~3208--3211). Not a contradiction --- but a convention, not a derived
compatibility. \textbf{Verdict: partially delivers (I): a convention is documented;
derived compatibility is unserved. Q high / K high / S see
Section~\ref{sec:notices}.} This is the current F-R2b/C6 classification,
not an E+ compatibility finding.

\subsection{F-R4a --- the \texorpdfstring{$j^2$}{j\texttwosuperior}-load: addressed and assigned}\label{sec:r4a}

\emph{Minimal predicate: the load is addressed --- its origin, propagation,
and compensating site are named; it does not vanish unbooked in one
coordinate and return silently.} The load arises only at the odd semistable
places (Thm.~4.2.2.1, eqs.~(4.2.2.2) vs.~(4.2.2.3); extract ll.~1544--1570);
it survives into the numerical evaluation (the $\sum_j j^2/\ell^{*2}$ of
\S\,9.9; extract ll.~6476--6498); and its return is assigned to a forced
renormalization at the complementary places, with the product formula as the
booking instrument (Cor.~4.2.2.4, extract ll.~1600--1614; Thm.~4.6.1(4),
extract ll.~1786--1791). The feared scenario --- silent vanishing and
unbooked return --- does not occur in the text.
\textbf{Verdict: delivers (row predicate). Q high / K high / S not checked.}
The historical address predicate does not quantify all loads, signs,
residuals, and compensating budget entries required by current C5;
no current C5/E+ finding follows. A recorded citation inconsistency: the source refers to the same relation
alternately as Theorem~4.2.2.1(3) and~(4); the referent is unambiguous.

\subsection{F-R4b --- quantitative balance: open}\label{sec:r4b}

\emph{Minimal predicate: the compensation is quantitatively booked ---
amount, uniqueness, and budget split shown.} The existence of the compensating
renormalization is inferred from the product-formula constraint in a one-line
proof (``This is clear from Theorem~4.2.2.1 and the global constraint placed
by the requirement that the product formula holds''; extract ll.~1613--1614).
The residual load is \emph{located but not quantified}; contract component~5
(separated indeterminacy budget) is served qualitatively, not quantitatively.
\textbf{Verdict: partially delivers. Q high / K high / S not checked.}

\subsection{F-R6b --- order relation, form level: written-out proof with an
addressable core step}\label{sec:r6f}

\emph{Minimal predicate: the order relation is stated as a theorem with a
written-out proof whose steps are addressable.} ATS~III, Theorem~9.11.1
(extract ll.~6794--6806; in Mochizuki's notation Cor.~9.11.1.1, extract
ll.~6882--6888) states the volume lower bound; the evaluation exists in the
container as an adelic log-volume computed as the sum of local contributions
(``one can compute log Vol (\dots) by taking the sum \dots of the local
contributions at all primes''; ATS~IV, \S\,1.5, extract ll.~764--769); this is
the contract component~6 evaluation audited in Section~\ref{sec:r6cc}, not
part of the order-relation finding recorded here. The cited mechanism is source-specific: ATS~III,
Lemma~9.10.7.1 (PDF p.~125) treats a superset $S'$ of a tensor-product
box, with integral $\tau_i\in O_{E_i}$, constituents $s_i\in\tau_i O_{E_i}$,
and weights $0<\gamma_i\le1$. Theorem~9.11.1 (PDF p.~127)
uses the asserted inclusion of that box in the construction.
Neither statement is the generic proposition that the volume of
an arbitrary set bounds the absolute value of each element.
Our v0.2 paraphrase of that kind was false: with normalized
additive Haar measure $\mu(\mathbb Z_p)=1$, the translated ball
\[
 U=1+p\mathbb Z_p,\qquad 1\in U,\qquad
 \mu(U)=p^{-1}<|1|_p=1
\]
is a counterexample. It refutes our unrestricted paraphrase,
not Joshi's tensor-box lemma, Theorem~9.11.1, IUT, or abc.
In particular, the centered-module hypothesis in ATS~III,
Lemma~9.10.2.2 (PDF p.~123) cannot be silently dropped.

The historical form finding is limited to a stated theorem and an
addressable proof. Establishing that the actual constructed
theta-set supplies the box, integral structure, and applicable
weight hypotheses remains the separate A2 obligation. Named
proof addresses alone do not discharge current C7: each
load-bearing hypothesis, including the inclusion hypothesis,
needs its own applicable anchor. The availability of adelic
evaluation alone does not establish C6 compatibility.
\textbf{Verdict: delivers the form predicate. Q high / K high / S externally
marked open (A2).}

\subsection{F-R6a --- cross-container compatibility: set, and global-only by
the source's own restriction}\label{sec:r6cc}

\emph{Minimal predicate (contract component~6, global form): the connecting
compatibility between the two evaluation routes is derived; contract
component~7 in turn uses it to conclude $\nu_q(\mathrm{Obj}_q) \in H_\Theta$
under admissible linkage.} For the \emph{abc} conclusion the lower bound is computed
with $\varphi(y_0)$, the upper with $y_0$; they are connected by the middle
equality of ATS~IV, Theorem~6.10.1, whose complete proof reads: ``The middle
equality is a tautology using the fact that the number fields provided by
$y_0'$ and $y_0$ are identically normalized'' (T-d; extract ll.~3272--3273).
Two clarifications keep the standard honest. First, the certificate's order
clause is \emph{global}; the source's own restriction --- one ``cannot claim
that $z_p \le z_p'$ holds for all primes $p$ without violating the middle
equality'' (ATS~IV, \S\,1.5, extract ll.~816--819), and ``one cannot make local
term by term comparisons at each prime'' (extract ll.~772--775) --- is a
\emph{scope limit of the claim}, not a violated criterion. Second, what
remains unserved at the form level is the \emph{derivation} of the
cross-container compatibility: it is set by definitional normalization, and
the remaining obligations include the indexed E2 logarithmic-volume
comparison, ambient and integral structures, theta-set and hull
compatibility, and measure/weight conventions
(Section~\ref{sec:notices}, notice A1). The valuation-level normalization
identification does not by itself carry those obligations.
\textbf{Verdict: partially delivers (declared convention; derivation
unserved). Q high / K high / S see notices.}

\subsection{F-R8 --- the second mode: visible and declared}\label{sec:r8}

\emph{Minimal predicate: the second mode used by the bridge is visible and
its role declared (not: a complete inventory of all auxiliary modes).} The
second mode is the global Frobenius/log shift, the pair
$(y_0, y_0' = \varphi(y_0))$. It is named; its effect is directionally
characterized (absolute values of Tate parameters rise under $\varphi$;
ATS~IV, \S\,1.4(2), extract ll.~567--569); its asymmetric use is declared
(``the lower bound will be computed using the arithmeticoid
$y_0'=\varphi(y_0)$ while the upper bound will be computed using $y_0$ (the
upper bound is not independent of this choice)''; \S\,1.5, extract
ll.~733--743); it is geometrically interpreted (slope comparison of
Frobenius-shifted bundles; Rem.~6.10.3, extract ll.~3213--3235); and it is
flagged as a correction of the paper's own earlier release (``[The previous
release of this paper had omitted this Frobenius-shifting.]''; extract
l.~755). \textbf{Verdict: delivers (visibility predicate). Q high / K high /
S not checked.}

Mode inventory (visibility status only; completeness is \emph{not} claimed):
the Frobenius/log shift (visible, used); the countably infinite choice of the
base point $y_0$ (visible, unbooked; ATS~III, Rem.~4.4.2, extract
ll.~1723--1731); the intended-but-unavailable average over arithmeticoids
(visible, unavailable; ATS~IV, \S\,4.5(4)/(5): ``one is far from being able to
do this'', extract ll.~2100--2121); the auxiliary per-place quantities
$z_p, z_p'$ constrained only through their sums (visible, unbooked; ATS~IV
eqs.~(1.5.4)/(1.5.5), extract ll.~798--819).

%=====================================================================
\section{External Status Notices and Their Corrected Scope}\label{sec:notices}

The historical audit was run on 12~August 2026; notices A1--A3
were recorded on 13/15~August. Their former terminal interpretations
are corrected here as of 2~October 2026. These notices are source-reported
follow-up assessments, not proofs established by this form audit.
Parts~5 and~6 are now published working papers
\cite{part5,part6}; public availability is not a mathematical
acceptance or a promotion of their claims.

\begin{table}[htbp]
\centering\footnotesize
\begin{tabularx}{\textwidth}{@{}l >{\raggedright\arraybackslash}X >{\raggedright\arraybackslash}X@{}}
\toprule
Notice & Corrected scope & Effect on this audit \\
\midrule
A1 (13 Aug) & valuation-level normalization; indexed hull/measure/weight and E2 compatibility remain open & F-R2b and F-R6a remain partial \\
A2 (13 Aug) & ATS III tensor-box inclusion and applicable hypotheses remain open & historical proof-address finding retained; no C7 substance certification \\
A3 (15 Aug) & conditional effective-weight classification under fixed field/place and identity hypotheses & no unconditional F-R6a or global branch closure \\
\bottomrule
\end{tabularx}
\caption{Corrected reading of the historical notices, 2~October 2026.
Dates identify their origin; the current text limits their import.}
\label{tab:status-notices}
\end{table}

A1 concerns the middle equality in ATS~IV, Theorem~6.10.1
(PDF pp.~66--67), with the scope limits of \S\,1.5 and
Remarks~6.10.2--3. The normalization supplement's corrective
v1.2~\cite{gnc} distinguishes a valuation-level identification
from the further choices needed to compare indexed E2
logarithmic volumes. Declared normalization does not settle
ambient/integral structure, theta-set, hull, measure, weight,
or ordering compatibility. A suggested Frobenius mechanism
for set stability remains source-reported; it is not an
established E2 transport theorem. F-NEU-1 remains open.

A2 concerns ATS~III, Lemma~9.10.7.1 and Theorem~9.11.1
(PDF pp.~125 and~127). Its box-inclusion obligation is separate
from A1's ATS~IV normalization/measure obligation. Correcting our
generic volume paraphrase does not discharge A2 and does not
refute the source-specific lemma. The currently cited
Part~5 assessment~\cite{part5} is retained as a dated report,
not independently certified in this paper.

A3 is limited to quantities demonstrated to factor through
effective product-formula weights. The reported classification
requires a fixed number field and matched full set of places,
real effective weights, and an identity for all nonzero elements
under the stated absolute-value conventions; source-specific
pointwise normalization identities impose further conditions.
The full product formula alone leaves a common scalar; a
pointwise normalization identity is a stronger hypothesis.
Neither conclusion automatically applies to logarithmic volumes,
measures, theta-sets, or hulls. The all-place hypothesis cannot
silently be substituted for a non-archimedean-only place range.
Parts~5 and~6~\cite{part5,part6} require their own source-bound
review; no unconditional cross-container verdict, terminal
closure of reading branch (c), or positive E2 theorem is
imported here.

%=====================================================================
\section{Synthesis (Conservative)}\label{sec:synthesis}

Table~\ref{tab:synthesis} follows the corrected notices and
distinguishes historical address findings from current
component obligations (Section~\ref{sec:crosswalk}).

\begin{table}[htbp]
\centering\footnotesize
\begin{tabularx}{\textwidth}{@{}l >{\raggedright\arraybackslash}X >{\raggedright\arraybackslash}X@{}}
\toprule
Finding & Observed form & Obligation not certified \\
\midrule
F-R2a / C4 & explicit $j^2$ transport law & C6 dependency; typed evaluation identity \\
F-R2b / C6 & declared convention; I & derived cross-container compatibility (A1) \\
F-R4a / C5 & load and site addressed (historical predicate) & complete quantified budget \\
F-R4b / C5 & qualitative location; partial & amount, split, residual and sign accounting \\
F-R6b / C7 & theorem and proof addresses (historical predicate) & applicable hypothesis anchors; actual box inclusion (A2) \\
F-R6a / C6 & evaluation available; global scope stated; partial & typed evaluation identity; indexed E2 hull/measure/weight compatibility (A1) \\
F-R8 / E8 & Frobenius mode visible (historical predicate) & complete coverage/exclusion of order-relevant modes \\
\bottomrule
\end{tabularx}
\caption{Conservative selected-rows synthesis; no current full-certificate
or proof-quality score. Q/K source/criterion confidence is distinct
from unchecked or open mathematical substance.}
\label{tab:synthesis}
\end{table}

The source supplies identifiable scaling, compensation-site, and
mode statements. F-R2b and F-R6a remain partial; F-R4b remains a
located but unquantified budget. Historical F-R4a, F-R6b and
F-R8 findings concern narrower predicates than current C5, C7
and E8. These observations permit inspection of the named
obligations, not a retrospective full-contract pass.
The translated-ball counterexample corrects this paper's own
paraphrase only. The A1 and A2 obligations remain distinct;
the conditional A3 classification does not close either.

\begin{table}[htbp]
\centering\footnotesize
\begin{tabularx}{\textwidth}{@{}l >{\raggedright\arraybackslash}X >{\raggedright\arraybackslash}X@{}}
\toprule
Level & Current scope & Remaining burden \\
\midrule
Normalization & valuation-level identification/convention & typed/indexed comparison not established \\
Measure / E2 & no general positive-measure classification imported & ambient, hull, measure, weight and ordering compatibility \\
Set stability & candidate Frobenius mechanism source-reported & no established indexed E2 set/hull transport \\
Box / A2 & specific tensor-box hypotheses and inclusion & actual construction and applicable hypothesis anchors \\
Effective weights / A3 & fixed-field/all-place conditional classification & applicability, place-range reconciliation; no measure consequence \\
\bottomrule
\end{tabularx}
\caption{Current burden map, 2~October 2026. The former
unconditional and terminal imports are withdrawn; open obligations
are not new refutations of ATS.}
\label{tab:burden}
\end{table}

%=====================================================================
\clearpage
\section{Claim Boundaries}\label{sec:claims}

This audit makes no claim about the truth of the \emph{abc} conjecture, the
correctness of IUT, the correctness of the ATS series, or the ultimate
repairability of the transport step. It judges publicly posted preprint text
against a pre-specified contract form, with verbatim anchors; an unserved row
predicate localizes a burden, it does not refute a theorem. The audit is
restricted to the four commissioned rows; no completeness of the mode
inventory, no full-certificate pass, and no aggregation of the row verdicts
into a proof-quality score is claimed. The substance-level analysis and repair classifications are the
subjects of the published working papers Parts~5 and~6
\cite{part5,part6}; this paper imports only explicitly limited,
dated source reports, not a fresh acceptance of their mathematics.

% Local, scoped adaptation of ai_disclosure_STANDARD_en, 2026-10-02.
\section*{AI Disclosure}
\noindent
AI systems were used for research discussion, literature and source work,
mathematical analysis, text, translation, and typesetting. Earlier
AI-assisted assessments remain subject to the source and claim limits
stated in this working paper.
\medskip
\noindent
For this corrective version, quality assurance includes separate
source-bound mathematical and bilingual delta reviews, checks of
the source versions and bibliographic identities, compilation,
and visual inspection of both PDFs. These checks concern the
specified corrective manuscript delta; they do not certify the
ATS proof chain. No new computational experiment or whole-series
proof verification is reported.


%=====================================================================
\begin{thebibliography}{99}
\small
\raggedright
\sloppy
\setlength{\itemsep}{2pt plus 0.5pt minus 0.5pt}

\bibitem{joshiI}
K.~Joshi,
\emph{Construction of Arithmetic Teichmüller Spaces I},
arXiv:2106.11452v4 [math.AG], 2021 (revised 24~February 2025).

\bibitem{joshiII}
K.~Joshi,
\emph{Construction of Arithmetic Teichmüller Spaces II: Proof of a Local
Prototype of Mochizuki's Corollary~3.12},
arXiv:2303.01662v3 [math.AG], 2023 (revised 24~February 2025).

\bibitem{joshiIIhalf}
K.~Joshi,
\emph{Construction of Arithmetic Teichmüller Spaces II$\frac{1}{2}$:
Deformations of Number Fields},
arXiv:2305.10398v12 [math.AG], 2023 (revised 24~February 2025).

\bibitem{joshiIII}
K.~Joshi,
\emph{Construction of Arithmetic Teichmüller Spaces III: A ``Rosetta Stone''
and a Proof of Mochizuki's Corollary~3.12},
arXiv:2401.13508v4 [math.AG], 2024 (revised 24~February 2025).

\bibitem{joshiIV}
K.~Joshi,
\emph{Construction of Arithmetic Teichmüller Spaces IV: Proof of the
abc-Conjecture},
arXiv:2403.10430v2 [math.AG], 2024 (revised 24~February 2025).

\bibitem{mochizukiIUT}
S.~Mochizuki,
\emph{Inter-universal Teichmüller Theory I: Construction of Hodge
Theaters}, Publ.\ Res.\ Inst.\ Math.\ Sci.\ \textbf{57} (2021), no.~1/2,
3--207;
\emph{II: Hodge--Arakelov-Theoretic Evaluation}, \textit{ibid.}, 209--401;
\emph{III: Canonical Splittings of the Log-Theta-Lattice}, \textit{ibid.}, 403--626;
\emph{IV: Log-Volume Computations and Set-Theoretic Foundations},
\textit{ibid.}, 627--723.

\bibitem{scholzestix}
P.~Scholze and J.~Stix,
\emph{Why abc is still a conjecture},
unpublished manuscript, version of 16~July 2018.
Available at
\url{https://www.math.uni-bonn.de/people/scholze/WhyABCisStillaConjecture.pdf}

\bibitem{part1}
L.~Geiger,
\emph{The Comparison in IUTchIII, Corollary 3.12: From Attempted Repair
to Structural Diagnosis},
Zenodo, 2026, version v6.4.
DOI: \nolinkurl{10.5281/zenodo.23081604}.

\bibitem{part2}
L.~Geiger,
\emph{The Bridge-Contract Criterion: A Form Audit for Transport Arguments,
with an Instantiation in the abc Literature},
Zenodo, 2026, version 0.2.
DOI: \nolinkurl{10.5281/zenodo.21960477}.

\bibitem{part3}
L.~Geiger,
\emph{The Bridge Certificate: A Line-Level Audit Instrument for
Inter-Theoretic Transfer Arguments, with Four Applications to the abc Literature},
Zenodo, 2026, version 0.3.
DOI: \nolinkurl{10.5281/zenodo.23090371}.

\bibitem{part5}
L.~Geiger,
\emph{The Joshi ATS Series under the Bridge-Certificate Standard:
A Targeted Substance Audit and an All-Place Rigidity Classification
of Effective Normalization Weights},
Zenodo, 2026, version 0.2.
DOI: \nolinkurl{10.5281/zenodo.21956399}.

\bibitem{part6}
L.~Geiger,
\emph{What the Product Formula Allows: A Rigidity Theorem and an
Effective-Weight Test for the ATS Height Subchain},
Zenodo, 2026, version 0.2.
DOI: \nolinkurl{10.5281/zenodo.21952487}.

\bibitem{gnc}
L.~Geiger,
\emph{Global Normalization in Joshi's Arithmetic Teichmüller Theory:
What the Identification Carries --- and What Carries the Theorem (DRAFT)},
Zenodo, 2026, version 1.2.
DOI: \nolinkurl{10.5281/zenodo.23092166}.

\bibitem{landscapeA}
L.~Geiger,
\emph{The abc Landscape: Obstructions, Failed Routes, and HCT Diagnostics
(DRAFT)},
Zenodo, 2026.
Concept DOI: \nolinkurl{10.5281/zenodo.21916900}
(cited version 1.1, 13~August 2026, DOI: \nolinkurl{10.5281/zenodo.21924338}).

\end{thebibliography}

\end{document}
